\documentclass{chsmc}
\chsmcsetup{
	year = 2012,
	part = 1,
	type = A,
	show-answers = false,
}
\begin{document}

\section{填空题：本大题共 8 小题，每小题 8 分，共 64 分.}

\begin{enumerate}
	\item 设 $P$ 是函数 $y = x + \dfrac{2}{x}$（$x > 0$）的图象上任意一点，
		过点 $P$ 分别向直线 $y = x$ 和 $y$ 轴作垂线，垂足分别为 $A$, $B$，
		则 $\overrightarrow{PA}\cdot\overrightarrow{PB}$ 的值是 \blankline
	\item 设 $\triangle ABC$ 的内角 $A$, $B$, $C$ 的对边分别为 $a$, $b$, $c$，
		且满足 $a\cos B - b\cos A = \dfrac{3}{5}c$，则 $\dfrac{\tan A}{\tan B}$ 的值是 \blankline
	\item 设 $x$, $y$, $z \in [0,1]$，则
		$M = \sqrt{|x-y|} + \sqrt{|y-z|} + \sqrt{|z-x|}$ 的最大值是 \blankline
	\item 抛物线 $y^2 = 2px$（$p > 0$）的焦点为 $F$，准线为 $l$，$A$, $B$ 是抛物线上的两个动点，
		且满足 $\angle AFB = \dfrac{\pi}{3}$. 设线段 $AB$ 的中点 $M$ 在 $l$ 上的投影为 $N$，
		则 $\dfrac{|MN|}{|AB|}$ 的最大值是 \blankline
	\item 设同底的两个正三棱锥 $P$-$ABC$ 和 $Q$-$ABC$ 内接于同一个球.
		若正三棱锥 $P$-$ABC$ 的侧面与底面所成的角为 $45^\circ$，
		则正三棱锥 $Q$-$ABC$ 的侧面与底面所成角的正切值是 \blankline
	\item 设 $f(x)$ 是定义在 $\mathbf{R}$ 上的奇函数，且当 $x \geqslant 0$ 时，$f(x) = x^2$.
		若对任意的 $x \in [a,a+2]$，不等式 $f(x+a) \geqslant 2f(x)$ 恒成立.
		则实数 $a$ 的取值范围是 \blankline
	\item 满足 $\dfrac{1}{4} < \sin \dfrac{\pi}{n} < \dfrac{1}{3}$ 的所有正整数 $n$ 的和是 \blankline
	\item 某情报站有 A, B, C, D 四种互不相同的密码，每周使用其中的一种密码，
		且每周都是从上周未使用的三种密码中等可能地随机选用一种.
		设第 $1$ 周使用 A 种密码，那么第 $7$ 周也使用 A 种密码的概率是 \blankline（用最简分数表示）
\end{enumerate}

\section{解答题：本大题共 3 个小题，满分 56 分. 解答应写出文字说明、证明过程或演算步骤.}

\begin{enumerate}[resume]
	\item (本题满分 16 分) 已知函数
		$f(x) = a\sin x - \dfrac{1}{2}\cos 2x + a - \dfrac{3}{a} + \dfrac{1}{2}$，
		$a \in \mathbf{R}$ 且 $a \neq 0$.
    \begin{enumerate}[(1)]
      \item 若对任意 $x \in \mathbf{R}$，都有 $f(x) \leqslant 0$，求 $a$ 的取值范围;
      \item 若 $a \geqslant 2$，且存在 $x \in \mathbf{R}$，使得 $f(x) \leqslant 0$，求 $a$ 的取值范围.
    \end{enumerate}
	\item (本题满分 20 分) 已知数列 $\{a_n\}$ 的各项均为非零实数，
		且对于任意的正整数 $n$，都有 $(a_1+a_2+\cdots+a_n)^2 = a_1^3+a_2^3+\cdots+a_n^3$.
    \begin{enumerate}[(1)]
      \item 当 $n = 3$ 时，求所有满足条件的三项组成的数列 $a_1$, $a_2$, $a_3$;
      \item 是否存在满足条件的无穷数列 $\{a_n\}$，使得 $a_{2013} = -2012$？若存在，求出这样的无穷数列的一个通项公式；若不存在，说明理由.
    \end{enumerate}
	\item (本题满分 20 分) 如图，在平面直角坐标系 $xOy$ 中，菱形 $ABCD$ 的边长为 $4$，
		且 $|OB| = |OD| = 6$.
    \begin{enumerate}[(1)]
      \item 求证：$|OA|\cdot|OC|$ 为定值;
      \item 当点 $A$ 在半圆 $M\colon (x-2)^2 + y^2 = 4$（$2 \leqslant x \leqslant 4$）上运动时，求点 $C$ 的轨迹.
    \end{enumerate}
\end{enumerate}

\end{document}
